% GENERATED FILE. Edit canonical lessons, then run npm run generate:books.
% canonical-source-hash: fnv1a64:35c5d4962ca1c0c7
% canonical-lessons: HI-C78-signs, HI-C78-notices, HI-C78-pehla-paath, HI-R85-devanagari-recall-a, HI-R85-devanagari-recall-b

\chapter{Reading — Signs, Notices, and a First Paragraph}
\label{ch:hi-pehla-paath}
\hlchaptermodality{85}

\begin{chapteropening}
\textbf{By the end of this chapter:} \emph{I can read Devanagari off a sign, a notice, and a short paragraph, without translating word by word.}

Seventy-eight words of Hindi in a row, and not one of them new — the first time the track asks for a whole paragraph at once.
\end{chapteropening}

\section[praveś · nikās · khulā · band · dukān …]{(reading) — six signs}

\label{lesson:HI-C78-signs}



\begin{quote}
Every word below, you have already been taught. Not one of them is new.

What is new is that nobody is going to tell you which one is coming.
\end{quote}

\begin{tcolorbox}[breakable,skin=enhanced,colback=orange!6,colframe=orange!45!black,boxrule=0.5pt,arc=1mm,left=6pt,right=6pt,top=4pt,bottom=4pt,fonttitle=\bfseries,title={Read this}]
\begin{quote}
\hi{प्रवेश} \hi{निकास} \hi{खुला} \hi{बंद} \hi{दुकान} \hi{स्टेशन}
\end{quote}

Read down the list once. Do not sound each letter out — look at the whole word and let it arrive.

Now again, and notice which ones came instantly and which ones you had to assemble. Both are fine. The difference is what practice removes.
\end{tcolorbox}

\begin{culture}
\textbf{\hi{प्रवेश}} and \textbf{\hi{निकास}} are the two you will meet at every door in India, and they are worth knowing by shape rather than by letter, because you usually read them while walking.

A sign is the one kind of Hindi that gives you no sentence to lean on — no verb at the end, no \textbf{\hi{है}} to tell you what is going on. One word, and it has to carry everything.

That is why this is the place reading starts. Six words is not much to hold, and holding them without help is the whole exercise.
\end{culture}

\subsection*{Your turn}
\begin{itemize}
\raggedright
  \item \emph{Say it:} the six words down the list, once, without stopping
\end{itemize}

\emph{Twice through:}

\begin{tcolorbox}[breakable,colback=teal!4,colframe=teal!35!black,title={Before you move on}]
How many of the six were new? (\textbf{None.})

What was new, then? (\textbf{Nobody said them first.} You read them.)
\end{tcolorbox}
\section[dukān band hai · sṭeśan khulā hai]{(reading) — six notices}

\label{lesson:HI-C78-notices}



\begin{quote}
The last lesson gave you single words. This one gives each word something to stand on.

Still nothing new — the same signs, with \textbf{\hi{है}} behind them.
\end{quote}

\begin{tcolorbox}[breakable,skin=enhanced,colback=orange!6,colframe=orange!45!black,boxrule=0.5pt,arc=1mm,left=6pt,right=6pt,top=4pt,bottom=4pt,fonttitle=\bfseries,title={Read this}]
\begin{quote}
\hi{दुकान} \hi{बंद} \hi{है।} \hi{स्टेशन} \hi{खुला} \hi{है।} \hi{यहाँ} \hi{पानी} \hi{है।} \hi{टिकट} \hi{यहाँ} \hi{है।} \hi{प्रवेश} \hi{यहाँ} \hi{है।} \hi{कृपया} \hi{धीरे।}
\end{quote}

Read down once, without stopping.

Again — and this time notice you did not have to work out what was shut. The notice said so.
\end{tcolorbox}

\begin{culture}
\textbf{\hi{दुकान} \hi{बंद} \hi{है।}} Three words, and the last one is the verb. Hindi puts it there every time, so the end of a notice is where you find out whether something \emph{is} or \emph{is not}.

That is worth more than it looks. It means you can read the first two words of any notice, guess, and then have your guess confirmed or corrected by the last word — which is exactly how a fluent reader reads.

The last line breaks the pattern: \textbf{\hi{कृपया} \hi{धीरे।}} No verb at all. Signs drop the verb when nobody could mistake the meaning, and \emph{please, slowly} needs no help.
\end{culture}

\subsection*{Your turn}
\begin{itemize}
\raggedright
  \item \emph{Say it:} the six notices, once, without stopping
\end{itemize}

\emph{Twice through:}

\begin{tcolorbox}[breakable,colback=teal!4,colframe=teal!35!black,title={Before you move on}]
Where does a Hindi notice keep its verb? (\textbf{At the end.})

And how many new words were in these six lines? (\textbf{None.})
\end{tcolorbox}
\section[merā nām aruṇ hai. maiṁ bhārat meṁ …]{(reading) — the first paragraph}

\label{lesson:HI-C78-pehla-paath}



\begin{quote}
Six words, then six lines. Now a whole paragraph.

Ninety-one words, and the rule has not changed: not one of them is new.
\end{quote}

\begin{tcolorbox}[breakable,skin=enhanced,colback=orange!6,colframe=orange!45!black,boxrule=0.5pt,arc=1mm,left=6pt,right=6pt,top=4pt,bottom=4pt,fonttitle=\bfseries,title={Read this}]
\begin{quote}
\hi{मेरा} \hi{नाम} \hi{अरुण} \hi{है।} \hi{मैं} \hi{भारत} \hi{में} \hi{रहता} \hi{हूँ।} \hi{मेरा} \hi{घर} \hi{छोटा} \hi{है}, \hi{लेकिन} \hi{अच्छा} \hi{है।} \hi{घर} \hi{में} \hi{एक} \hi{कमरा}, \hi{एक} \hi{दरवाज़ा} \hi{और} \hi{दो} \hi{खिड़की} \hi{हैं।} \hi{सुबह} \hi{मैं} \hi{चाय} \hi{और} \hi{दूध} \hi{पीता} \hi{हूँ।} \hi{फिर} \hi{मैं} \hi{रोटी} \hi{खाता} \hi{हूँ।} \hi{आज} \hi{मैं} \hi{बाज़ार} \hi{जाता} \hi{हूँ}, \hi{क्योंकि} \hi{मुझे} \hi{फल} \hi{पसंद} \hi{हैं।} \hi{बाज़ार} \hi{में} \hi{एक} \hi{दुकान} \hi{है।} \hi{वहाँ} \hi{एक} \hi{आदमी} \hi{काम} \hi{करता} \hi{है।} \hi{वह} \hi{मेरा} \hi{दोस्त} \hi{है।} \hi{वह} \hi{हिंदी} \hi{बोलता} \hi{है।} \hi{वह} \hi{किताब} \hi{पढ़ता} \hi{है}, \hi{लेकिन} \hi{वह} \hi{रात} \hi{में} \hi{दुकान} \hi{में} \hi{काम} \hi{करता} \hi{है।} \hi{फिर} \hi{मैं} \hi{घर} \hi{आता} \hi{हूँ।} \hi{रात} \hi{में} \hi{मैं} \hi{किताब} \hi{पढ़ता} \hi{हूँ।}
\end{quote}

Read once without stopping. Let the sentences arrive without translation.

Now read it again, and notice how little work it took.
\end{tcolorbox}

\begin{culture}
Look at what is holding it together, because it is not vocabulary.

\textbf{\hi{और}} adds. \textbf{\hi{लेकिन}} holds two true things apart. \textbf{\hi{क्योंकि}} gives a reason. \textbf{\hi{फिर}} puts one thing after another. Each of those was a chapter, and each was small.

This is the first time the track has asked you to hold a whole paragraph at once — and the reason it can is that everything in it was paid for, one chapter at a time.
\end{culture}

\subsection*{Your turn}
\begin{itemize}
\raggedright
  \item \emph{Say it:} the paragraph aloud, once, without stopping
\end{itemize}

\emph{Twice through:}

\begin{tcolorbox}[breakable,colback=teal!4,colframe=teal!35!black,title={Before you move on}]
How many words in that paragraph were new to you? (\textbf{None.})

And what made it readable? (Everything in it had already been \textbf{paid for}, one chapter at a time.)
\end{tcolorbox}
\section[pahlī gyārah]{(\hi{म} \hi{र} \hi{न} \hi{अ} \hi{प} \hi{ह} \hi{य} \hi{त} \hi{्} \hi{ा} \hi{ी}) — eleven, and three kinds of thing}

\label{lesson:HI-R85-devanagari-recall-a}



\begin{quote}
Two letters from the book's first script pages. Read \textbf{\hi{म}}, then \textbf{\hi{न}}, and name each.
\end{quote}

\subsection*{Script: eight letters, one killer, two hangers}
They arrived one or two per chapter and no page has gathered them. Sorted by what kind of thing each one is:

\noindent
\begin{tabularx}{\linewidth}{@{}>{\raggedright\arraybackslash}X>{\raggedright\arraybackslash}X@{}}
\toprule
\textbf{kind} & \textbf{pieces} \\
\midrule
\textbf{consonants} — each carries a built-in \emph{a} & \textbf{\hi{म}} \emph{ma} · \textbf{\hi{र}} \emph{ra} · \textbf{\hi{न}} \emph{na} · \textbf{\hi{प}} \emph{pa} · \textbf{\hi{ह}} \emph{ha} · \textbf{\hi{य}} \emph{ya} · \textbf{\hi{त}} \emph{ta} \\
\textbf{a standing vowel} & \textbf{\hi{अ}} \emph{a} — the full-size letter, for when a word opens on the vowel \\
\textbf{the killer} & \textbf{\hi{्}} — removes the built-in \emph{a} \\
\textbf{hangers} & \textbf{\hi{ा}} \emph{ā} · \textbf{\hi{ी}} \emph{ī} — replace the built-in \emph{a} \\
\bottomrule
\end{tabularx}

\textbf{\hi{अ}} is the one worth pausing on, because it is the same sound the consonants already carry. A consonant has \emph{a} inside it; \textbf{\hi{अ}} is what you write when there is no consonant to put it in.

\subsection*{Your turn}
The three marks, cold, and what each does to the letter it touches.

Read these aloud:

\begin{itemize}
\raggedright
  \item \textbf{\hi{्}} — say what it takes away
  \item \textbf{\hi{ा}} — say what it replaces, and with what
  \item \textbf{\hi{ी}} — the same, and say how it differs from the one above
\end{itemize}

All three depend on one fact: a bare consonant is not silent. It already has a vowel, and these either remove it or swap it.

\begin{tcolorbox}[breakable,colback=teal!4,colframe=teal!35!black,title={Before you move on}]
What vowel is inside a bare Devanagari consonant? (\textbf{A short \emph{a}.}) What does \textbf{\hi{्}} do? (\textbf{Removes it.}) And when do you need \textbf{\hi{अ}} rather than relying on that built-in vowel? (\textbf{When the word opens on the vowel.})
\end{tcolorbox}
\section[dūsrī gyārah]{(\hi{स} \hi{ष} \hi{ल} \hi{ठ} \hi{ण} \hi{ऋ} \hi{ो} \hi{ौ} \hi{औ} \hi{ज़}) — the thin eleven}

\label{lesson:HI-R85-devanagari-recall-b}



\begin{quote}
Read \textbf{\hi{स}}, then \textbf{\hi{ल}}. Two consonants you have been reading for a long time without being asked for them by name.
\end{quote}

\subsection*{Script: five consonants, and the retroflex row they belong to}
\noindent
\begin{tabularx}{\linewidth}{@{}>{\raggedright\arraybackslash}X>{\raggedright\arraybackslash}X>{\raggedright\arraybackslash}X@{}}
\toprule
\textbf{letter} & \textbf{sound} & \textbf{what marks it} \\
\midrule
\textbf{\hi{स}} & \emph{sa} & the plain sibilant \\
\textbf{\hi{ल}} & \emph{la} & the plain lateral \\
\textbf{\hi{ष}} & \emph{ṣa} & \textbf{retroflex} — tongue curled back \\
\textbf{\hi{ठ}} & \emph{ṭha} & \textbf{retroflex}, with breath behind it \\
\textbf{\hi{ण}} & \emph{ṇa} & \textbf{retroflex} nasal \\
\bottomrule
\end{tabularx}

Three of the five are retroflex, and that is the group worth seeing together: Hindi distinguishes a whole row of consonants made with the tongue curled back against the roof of the mouth from the dental row made against the teeth. The difference is inaudible to most English speakers at first and load-bearing in Hindi.

\subsection*{Your turn: a vowel written like a letter, and a dot}
\begin{itemize}
\raggedright
  \item \emph{Read it:} \textbf{\hi{ऋ}} — inherited from Sanskrit, a vowel that behaves like a consonant. Hindi says it \emph{ri}
  \item \emph{Write it:} \textbf{\hi{ऋ}} once, from memory
  \item \emph{Read it:} \textbf{\hi{औ}} and \textbf{\hi{ौ}} — the same vowel standing alone, then hanging off a consonant
  \item \emph{Read it:} \textbf{\hi{ो}} — the hanger for \emph{o}
\end{itemize}

\textbf{\hi{औ}} against \textbf{\hi{ौ}} is the pairing the whole system runs on: a vowel gets a full-size letter when nothing precedes it and a hanging sign when a consonant does.

\begin{tcolorbox}[breakable,colback=teal!4,colframe=teal!35!black,title={Before you move on}]
What does the dot under \textbf{\hi{ज़}}, \textbf{\hi{फ़}}, \textbf{\hi{ख़}}, \textbf{\hi{क़}} tell you? (\textbf{The sound is borrowed} — Persian or Arabic, and not native to Devanagari.) Which of the eleven is a vowel that behaves like a consonant? (\textbf{\hi{ऋ}}.) Name the retroflex nasal. (\textbf{\hi{ण}}.)
\end{tcolorbox}
